\appendix
\section{Apéndice: Thue-Morse}

\begin{frame}
  \vspace{1.5cm}
  \centering
  {\LARGE Apéndice: Thue-Morse}
\end{frame}

\begin{frame}{Thue-Morse}
  \(\tau_0 = 0\).
  Cada término es el anterior, seguido de su complemento.

  \vspace{0.7em}
  \begin{center}
  \renewcommand{\arraystretch}{1.35}
  \begin{tabular}{@{}r@{\quad}l@{}}
    \(\tau_1\) & \textcolor{shared}{\texttt{0}}\hspace{0.45em}\textcolor{tail}{\texttt{1}} \\
    \(\tau_2\) & \textcolor{shared}{\texttt{01}}\hspace{0.45em}\textcolor{tail}{\texttt{10}} \\
    \(\tau_3\) & \textcolor{shared}{\texttt{0110}}\hspace{0.45em}\textcolor{tail}{\texttt{1001}} \\
    \(\tau_4\) & \textcolor{shared}{\texttt{01101001}}\hspace{0.45em}\textcolor{tail}{\texttt{10010110}}
  \end{tabular}
  \end{center}

  \vspace{0.6em}
  {\color{shared}\rule{0.7em}{0.7em}} la primera mitad, que es \(\tau_{i-1}\)
  \qquad
  {\color{tail}\rule{0.7em}{0.7em}} su complemento
\end{frame}

\begin{frame}{Por qué colisionan}
  \(W\) es el hash polinomial sin tomar módulo.
  La segunda mitad de \(\tau_n\) es el complemento de la primera, así que
  \[
    W(\overline{\tau_n}) - W(\tau_n)
      = \prod_{i=0}^{n-1} \bigl(1 - b^{2^{i}}\bigr).
  \]

  Si \(b^{2^{k}} \equiv 1 \pmod{M}\) para algún \(k < n\), ese factor es \(0\) módulo \(M\).
  Los dos hashes coinciden. El módulo puede ser cualquiera.

  \medskip
  Si \(M = 2^{64}\) y \(b\) es impar, cada factor aporta potencias de \(2\):
  \(2^{i+1}\) divide a \(1 - b^{2^{i}}\), y el producto junta al menos \(n(n+1)/2\) factores \(2\).
  Con \(n = 11\) eso pasa de \(64\).
  \(\tau_{11}\) y su complemento, de largo \(2048\), colisionan para cualquier base impar.
\end{frame}

\begin{frame}{Un ejemplo}
  Base \(b = 3\), módulo \(2^{10} = 1024\), con \(n = 4\).

  \begin{center}
  \texttt{0110 1001 1001 0110}\\[0.25em]
  \texttt{1001 0110 0110 1001}\\[0.7em]
  Los dos hashes valen \(464\).
  \end{center}

  \vspace{0.45em}
  Otro módulo, que no es potencia de \(2\): base \(b = 4\), módulo \(15\).
  Como \(4^{2} \equiv 1 \pmod{15}\), \(\tau_2 = \texttt{0110}\) y su complemento \(\texttt{1001}\) tienen hash \(5\).
\end{frame}

\begin{frame}[fragile]{Cómo generarlas}
\begin{lstlisting}
vector<int> thue(int n) {
    vector<int> s = {0};
    for (int k = 0; k < n; k++) {
        int m = (int) s.size();
        for (int i = 0; i < m; i++)
            s.push_back(1 - s[i]);
    }
    return s;
}

uint64_t H(vector<int> const& s, uint64_t b, uint64_t mod) {
    uint64_t h = 0, p = 1;
    for (int x : s) {
        h = (h + (uint64_t) x * p) % mod;
        p = p * b % mod;
    }
    return h;
}
// H(thue(4), 3, 1024) y el hash del complemento valen 464
\end{lstlisting}
\end{frame}
